\documentclass[10pt,a4paper,openany]{book}
\makeatletter
\def\input@path{{Style/}}
\makeatother
\usepackage{TLPortableCours}

\begin{document}
\TLPortableTitlePage{5}{Introduction et modélisation des systèmes asservis}{Boucles, performances, modèles temporels et fréquentiels}{Version propre et modulaire issue des cours originaux}
\tableofcontents
\clearpage
\chapter{Introduction et modélisation des systèmes asservis}
\markboth{Systèmes asservis}{Systèmes asservis}

\begin{savoirs}
\begin{itemize}
  \item structure d'un système automatisé en boucle ouverte ou fermée ;
  \item critères de stabilité, rapidité, précision et dépassement ;
  \item transformée de Laplace et fonction de transfert ;
  \item modèles canoniques du premier et du second ordre ;
  \item FTBO, FTBF, schémas-blocs, perturbations et réponse fréquentielle.
\end{itemize}
\end{savoirs}

\begin{savoirfaire}
\begin{itemize}
  \item établir une fonction de transfert à partir d'équations physiques ;
  \item réduire un schéma-blocs et calculer une FTBF ;
  \item identifier un modèle à partir d'une réponse temporelle ;
  \item déterminer erreur statique, temps de réponse et dépassement ;
  \item tracer et exploiter un diagramme de Bode.
\end{itemize}
\end{savoirfaire}

% ==========================================================================
% COURS 05 -- PARTIE A : INTRODUCTION AUX SYSTÈMES ASSERVIS
% ==========================================================================

\section{Mise en situation}\label{sec:sa-mise-situation}

Un grand nombre de systèmes techniques doivent maintenir une grandeur autour d'une valeur souhaitée malgré les variations de fonctionnement et les perturbations extérieures. Une étuve doit conserver une température donnée, un véhicule doit respecter une vitesse de consigne, un robot doit rejoindre une position imposée et une nacelle stabilisée doit maintenir son orientation.

Dans tous ces cas, la commande ne peut pas être déterminée une fois pour toutes. Le système doit observer son propre comportement, comparer le résultat à l'objectif et corriger son action.

\begin{SAapplication}[title={Trois situations typiques}]
\begin{itemize}
  \item \textbf{Étuve thermique :} la puissance de chauffe est adaptée à l'écart entre température demandée et température mesurée.
  \item \textbf{Régulateur de vitesse :} le couple moteur varie lorsque la pente, le vent ou la charge du véhicule changent.
  \item \textbf{Robot de chirurgie :} la position réelle de l'organe terminal est mesurée pour garantir précision et sécurité.
\end{itemize}
\end{SAapplication}

\SASchemaEtuve

\section{Systèmes automatisés et systèmes bouclés}\label{sec:sa-definitions}

\subsection{Système automatisé}

\begin{SAdefinition}
Un \textbf{système automatisé} réalise une suite d'opérations sans intervention humaine permanente. Il comporte généralement une partie commande, une partie opérative, des interfaces d'entrée et des capteurs.
\end{SAdefinition}

La partie commande élabore les ordres. La partie opérative transforme l'énergie et agit sur la matière d'œuvre. Les capteurs transmettent des informations relatives à l'état du système ou de son environnement.

\SAChaineFonctionnelle

Un automatisme séquentiel peut fonctionner sans mesurer précisément la grandeur obtenue. À l'inverse, un système asservi repose sur une comparaison continue entre une consigne et une mesure.

\subsection{Boucle ouverte et boucle fermée}

\begin{SAdefinition}[title={Commande en boucle ouverte}]
Dans une commande en \textbf{boucle ouverte}, l'action de commande ne dépend pas de la sortie effectivement obtenue. Une perturbation ou une variation du procédé n'est donc pas corrigée automatiquement.
\end{SAdefinition}

Exemple : alimenter un moteur pendant une durée prédéterminée pour déplacer un chariot. La distance obtenue dépend alors de la charge, des frottements et de la tension d'alimentation.

\begin{SAdefinition}[title={Commande en boucle fermée}]
Dans une commande en \textbf{boucle fermée}, la sortie est mesurée puis comparée à la consigne. L'écart, appelé \textbf{erreur}, est utilisé pour élaborer la commande.
\end{SAdefinition}

\SASchemaBoucleGenerale

On note généralement :
\[
\varepsilon(t)=e(t)-m(t),
\]
où $e(t)$ est la consigne et $m(t)$ la mesure ramenée au comparateur.

\begin{SAattention}
La mesure et la consigne doivent être de même nature et exprimées avec la même échelle au niveau du comparateur. Un capteur et son conditionneur peuvent donc être nécessaires pour adapter la grandeur physique mesurée.
\end{SAattention}

\subsection{Asservissement et régulation}

\begin{itemize}
  \item Un \textbf{asservissement} vise à faire suivre une consigne variable à la sortie : position d'un axe, trajectoire d'un robot, vitesse d'un véhicule.
  \item Une \textbf{régulation} vise à maintenir une grandeur autour d'une valeur constante malgré les perturbations : température, pression, niveau, tension électrique.
\end{itemize}

Dans les deux cas, la structure mathématique est identique. La distinction porte principalement sur la nature de la consigne et sur l'objectif d'utilisation.

\section{Structure générale d'un système asservi}\label{sec:sa-structure}

\SAChaineAsservie

Les principaux constituants sont :

\begin{description}
  \item[Comparateur :] calcule l'erreur entre consigne et mesure.
  \item[Correcteur :] transforme l'erreur en commande afin d'améliorer les performances.
  \item[Procédé :] système physique commandé, parfois associé à un préactionneur et à un actionneur.
  \item[Capteur :] mesure la sortie ou une grandeur qui lui est liée.
  \item[Chaîne de retour :] adapte la mesure avant son retour au comparateur.
\end{description}

La perturbation agit sur le procédé sans être choisie par la commande. Elle peut être une force extérieure, un couple résistant, une variation de température ambiante ou une fluctuation de tension.

\TLSchemaPerturbe

\subsection{Grandeurs usuelles}

Dans le domaine de Laplace, on note souvent :
\[
E(p)\text{ la consigne},\qquad S(p)\text{ la sortie},\qquad
\varepsilon(p)\text{ l'erreur},\qquad Q(p)\text{ la perturbation}.
\]

Pour une chaîne directe $D(p)$ et une chaîne de retour $R(p)$ :
\[
\varepsilon(p)=E(p)-R(p)S(p),\qquad S(p)=D(p)\varepsilon(p).
\]

On obtient alors :
\[
\boxed{\frac{S(p)}{E(p)}=\frac{D(p)}{1+D(p)R(p)}}.
\]

Le produit $D(p)R(p)$ est la \textbf{fonction de transfert en boucle ouverte}. Le quotient précédent est la \textbf{fonction de transfert en boucle fermée}.

\section{Performances attendues}\label{sec:sa-performances}

Un système asservi n'est pas seulement jugé sur sa capacité à fonctionner. Il doit respecter des critères quantifiés issus du cahier des charges.

\SAReperePerformances

\subsection{Stabilité}

\begin{SAdefinition}
Un système est \textbf{stable} si une entrée bornée conduit à une sortie bornée et si, après une perturbation transitoire, la réponse revient vers un régime permanent.
\end{SAdefinition}

Une réponse qui diverge ou dont les oscillations augmentent est instable. Une réponse dont les oscillations persistent sans décroître est à la limite de stabilité.

\subsection{Rapidité}

La rapidité est caractérisée par un temps de réponse. Pour une réponse à un échelon, on utilise fréquemment le \textbf{temps de réponse à 5\%}, noté $t_{5\%}$.

\begin{SAdefinition}
Le temps de réponse à 5\% est le premier instant à partir duquel la sortie reste définitivement dans le tube compris entre $95\%$ et $105\%$ de sa valeur finale.
\end{SAdefinition}

\begin{SAattention}
Le tube à 5\% est construit autour de la \textbf{valeur finale de la sortie}, et non autour de la valeur de la consigne lorsque le système présente une erreur statique.
\end{SAattention}

\subsection{Précision}

La précision traduit l'écart entre la sortie et la consigne en régime permanent. On définit l'erreur statique :
\[
\varepsilon_s=\lim_{t\to+\infty}\varepsilon(t).
\]

Pour une consigne constante, un système est parfaitement précis si $\varepsilon_s=0$.

\subsection{Dépassement}

Le dépassement relatif du premier maximum est défini par :
\[
D_1=100\,\frac{s_{\max}-s_\infty}{s_\infty}\quad(\%).
\]

Un dépassement important peut être incompatible avec la sécurité ou avec les contraintes mécaniques. Il n'existe pas pour un système du premier ordre idéal, mais apparaît fréquemment pour les systèmes du second ordre sous-amortis.

\subsection{Robustesse}

La robustesse est la capacité à conserver des performances acceptables lorsque le modèle varie : masse transportée, frottements, température, vieillissement ou imprécision des paramètres.

\begin{SAmethode}[title={Lecture d'un cahier des charges}]
Pour chaque exigence, identifier :
\begin{enumerate}
  \item la grandeur observée et son unité ;
  \item le scénario d'essai et le signal d'entrée ;
  \item le critère : stabilité, temps de réponse, erreur, dépassement ;
  \item la valeur limite et la tolérance.
\end{enumerate}
\end{SAmethode}

\section{Signaux tests}\label{sec:sa-signaux-tests}

Pour comparer les systèmes, on les sollicite avec des entrées normalisées.

\SASignauxTests

\begin{itemize}
  \item L'\textbf{impulsion} met en évidence la dynamique propre du système.
  \item L'\textbf{échelon} représente un changement brusque de consigne et permet d'étudier rapidité, dépassement et erreur statique.
  \item La \textbf{rampe} représente une consigne évoluant à vitesse constante et permet d'étudier l'erreur de traînage.
\end{itemize}

Dans le domaine de Laplace, pour une amplitude $E_0$ :
\[
\delta(t)\longleftrightarrow 1,\qquad
E_0u(t)\longleftrightarrow \frac{E_0}{p},\qquad
E_0t\,u(t)\longleftrightarrow \frac{E_0}{p^2}.
\]

\begin{SAapplication}[title={Choisir le bon essai}]
Pour mesurer un temps de réponse et un dépassement, on impose un échelon. Pour évaluer le suivi d'une trajectoire à vitesse constante, on impose une rampe. Pour identifier rapidement une dynamique, un essai indiciel est souvent le plus simple à réaliser expérimentalement.
\end{SAapplication}

% ==========================================================================
% COURS 05 -- PARTIE B : MODÉLISATION DES SLCI
% ==========================================================================

\section{Modéliser un système continu linéaire invariant}\label{sec:sa-slci}

L'objectif de la modélisation est de relier les grandeurs d'entrée et de sortie par un modèle suffisamment simple pour prévoir le comportement, dimensionner les composants et choisir un correcteur.

\SACarteModelisation

\subsection{Hypothèses d'un SLCI}

Un système est dit \textbf{continu, linéaire et invariant} lorsque :

\begin{itemize}
  \item ses grandeurs évoluent continûment avec le temps ;
  \item le principe de superposition est applicable ;
  \item ses paramètres ne dépendent pas explicitement du temps.
\end{itemize}

La linéarité implique que, si les entrées $e_1$ et $e_2$ produisent les sorties $s_1$ et $s_2$, alors l'entrée $\alpha e_1+\beta e_2$ produit $\alpha s_1+\beta s_2$.

\begin{SAattention}
Les systèmes réels présentent souvent des non-linéarités : saturation, jeu, seuil, frottement sec, hystérésis ou limitation de course. Le modèle linéaire n'est valable que dans un domaine de fonctionnement précisé.
\end{SAattention}

\subsection{Modélisation de connaissance}

La modélisation de connaissance est obtenue à partir des lois physiques : lois de Newton, relations électriques, bilans thermiques, conservation de la matière ou relations hydrauliques.

\begin{SAapplication}[title={Moteur à courant continu simplifié}]
Pour un moteur sans frottement sec :
\[
\begin{aligned}
 u(t)&=R i(t)+L\frac{di(t)}{dt}+e(t),\\
 e(t)&=K_e\omega(t),\\
 J\frac{d\omega(t)}{dt}+f\omega(t)&=K_t i(t)-C_r(t).
\end{aligned}
\]
Ces équations couplent la partie électrique et la partie mécanique. Elles permettent d'établir une fonction de transfert entre tension, couple résistant et vitesse.
\end{SAapplication}

\subsection{Linéarisation autour d'un point de fonctionnement}

Pour une relation non linéaire $s=f(e)$, on choisit un point de fonctionnement $(e_0,s_0)$ et on étudie les petites variations :
\[
\Delta e=e-e_0,\qquad \Delta s=s-s_0.
\]
Au premier ordre :
\[
\Delta s\simeq f'(e_0)\Delta e.
\]
Le gain local $f'(e_0)$ dépend du point de fonctionnement. La validité du modèle doit donc être vérifiée sur le domaine d'utilisation.

\section{Transformée de Laplace}\label{sec:sa-laplace}

La transformée de Laplace remplace les équations différentielles par des équations algébriques en la variable complexe $p$.

\begin{SAdefinition}
Pour une fonction causale $f(t)$, sa transformée de Laplace est :
\[
F(p)=\mathcal{L}\{f(t)\}=\int_0^{+\infty}f(t)e^{-pt}\,dt.
\]
\end{SAdefinition}

\subsection{Propriétés utiles}

Sous l'hypothèse de conditions initiales nulles :
\[
\mathcal{L}\left\{\frac{df}{dt}\right\}=pF(p),\qquad
\mathcal{L}\left\{\frac{d^n f}{dt^n}\right\}=p^nF(p).
\]

Plus généralement :
\[
\mathcal{L}\left\{\frac{df}{dt}\right\}=pF(p)-f(0^+).
\]

Les transformées usuelles sont :
\[
1\longleftrightarrow \frac1p,\qquad
t\longleftrightarrow \frac1{p^2},\qquad
e^{-at}\longleftrightarrow \frac1{p+a},
\]
\[
\sin(\omega t)\longleftrightarrow \frac{\omega}{p^2+\omega^2},\qquad
\cos(\omega t)\longleftrightarrow \frac{p}{p^2+\omega^2}.
\]

\begin{SAmethode}[title={Passage d'une équation différentielle à Laplace}]
\begin{enumerate}
  \item Écrire l'équation différentielle en séparant entrée et sortie.
  \item Préciser les conditions initiales.
  \item Transformer chaque dérivée en puissance de $p$.
  \item Regrouper les termes en $S(p)$ et en $E(p)$.
  \item Former le rapport $S(p)/E(p)$ pour conditions initiales nulles.
\end{enumerate}
\end{SAmethode}

\section{Fonction de transfert}\label{sec:sa-fonction-transfert}

\begin{SAdefinition}
La \textbf{fonction de transfert} d'un SLCI est le rapport entre la transformée de Laplace de la sortie et celle de l'entrée, pour des conditions initiales nulles :
\[
H(p)=\frac{S(p)}{E(p)}.
\]
\end{SAdefinition}

Pour une équation différentielle générale :
\[
a_n s^{(n)}+\cdots+a_1\dot{s}+a_0s
=b_m e^{(m)}+\cdots+b_1\dot{e}+b_0e,
\]
on obtient :
\[
H(p)=\frac{b_mp^m+\cdots+b_1p+b_0}{a_np^n+\cdots+a_1p+a_0}.
\]

\subsection{Pôles, zéros, ordre et gain statique}

Les racines du numérateur sont les \textbf{zéros}. Les racines du dénominateur sont les \textbf{pôles}. L'ordre du système est le degré du dénominateur après simplification des facteurs communs.

Lorsque la limite existe :
\[
K=H(0)
\]
est le gain statique. Il représente le rapport sortie/entrée en régime permanent pour une entrée constante.

\begin{SAattention}
Une simplification algébrique entre un pôle et un zéro ne signifie pas nécessairement que la dynamique physique correspondante a disparu. Une compensation exacte est rarement robuste lorsque les paramètres sont incertains.
\end{SAattention}

\subsection{Théorèmes des valeurs initiale et finale}

Sous leurs conditions d'application :
\[
\lim_{t\to0^+}f(t)=\lim_{p\to+\infty}pF(p),
\]
\[
\lim_{t\to+\infty}f(t)=\lim_{p\to0}pF(p).
\]

Le théorème de la valeur finale nécessite notamment que les pôles de $pF(p)$ soient à partie réelle strictement négative, à l'exception éventuelle d'un pôle simple à l'origine.

\section{Modèles canoniques}\label{sec:sa-modeles-canoniques}

\subsection{Gain pur}

Un gain pur est défini par :
\[
H(p)=K,\qquad s(t)=K e(t).
\]
Il ne possède aucune dynamique et représente une relation proportionnelle instantanée dans le domaine de validité du modèle.

\subsection{Intégrateur}

\[
H(p)=\frac{K}{p}.
\]
Pour un échelon d'amplitude $E_0$, la sortie est une rampe :
\[
s(t)=K E_0 t.
\]
Un intégrateur accumule l'entrée et introduit une phase de $-90^\circ$ en régime sinusoïdal.

\subsection{Premier ordre}

La forme canonique est :
\[
\boxed{H(p)=\frac{K}{1+\tau p}},
\]
où $K$ est le gain statique et $\tau>0$ la constante de temps.

L'équation différentielle associée est :
\[
\tau\dot{s}(t)+s(t)=K e(t).
\]

Pour un échelon d'amplitude $E_0$ et une sortie initialement nulle :
\[
\boxed{s(t)=K E_0\left(1-e^{-t/\tau}\right)}.
\]

\SAReponsePremierOrdre

Valeurs caractéristiques :
\[
s(\tau)=0{,}632\,K E_0,\qquad
t_{5\%}\simeq3\tau,
\]
\[
t_{2\%}\simeq4\tau,\qquad
\left.\frac{ds}{dt}\right|_{0^+}=\frac{K E_0}{\tau}.
\]

\begin{SAmethode}[title={Identifier un premier ordre sur une réponse indicielle}]
\begin{enumerate}
  \item Déterminer la variation finale $\Delta s_\infty$ et en déduire $K=\Delta s_\infty/\Delta e$.
  \item Mesurer $\tau$ à l'instant où la réponse atteint $63{,}2\%$ de sa variation finale.
  \item Vérifier la cohérence avec $t_{5\%}\simeq3\tau$ ou avec la tangente à l'origine.
\end{enumerate}
\end{SAmethode}

\subsection{Second ordre}

La forme canonique est :
\[
\boxed{H(p)=\frac{K\omega_0^2}{p^2+2\xi\omega_0p+\omega_0^2}}
=\frac{K}{1+\dfrac{2\xi}{\omega_0}p+\dfrac{p^2}{\omega_0^2}}.
\]

$\omega_0$ est la pulsation propre non amortie et $\xi$ le coefficient d'amortissement.

\begin{itemize}
  \item $\xi>1$ : régime apériodique ;
  \item $\xi=1$ : régime critique ;
  \item $0<\xi<1$ : régime pseudo-périodique ;
  \item $\xi=0$ : oscillations non amorties.
\end{itemize}

Pour $0<\xi<1$, la pulsation amortie est :
\[
\omega_a=\omega_0\sqrt{1-\xi^2},
\]
et la pseudo-période :
\[
T_a=\frac{2\pi}{\omega_a}.
\]

\SAReponseSecondOrdre

Le premier dépassement relatif est :
\[
\boxed{D_1=100\exp\left(-\frac{\pi\xi}{\sqrt{1-\xi^2}}\right)}.
\]

Une approximation courante du temps de réponse à 5\% est :
\[
t_{5\%}\simeq\frac{3}{\xi\omega_0},
\]
à utiliser avec prudence lorsque $\xi$ est proche de 1.

\begin{SAapplication}[title={Lecture physique des paramètres}]
Une augmentation de $\omega_0$ accélère globalement la dynamique. Une augmentation de $\xi$ réduit les oscillations et le dépassement, mais un amortissement trop important peut ralentir la réponse. Le choix résulte donc d'un compromis.
\end{SAapplication}

\section{Modèles avec retard et réduction d'ordre}\label{sec:sa-retard-reduction}

Un retard pur $T$ est représenté par :
\[
H_r(p)=e^{-Tp}.
\]
Il ne modifie pas le gain mais ajoute une phase $-\omega T$. Il est souvent approché par une fraction rationnelle de Padé :
\[
e^{-Tp}\simeq\frac{1-\dfrac{T}{2}p}{1+\dfrac{T}{2}p}.
\]

Lorsque plusieurs constantes de temps sont très différentes, les dynamiques rapides peuvent parfois être négligées devant la dynamique dominante. Une réduction d'ordre n'est acceptable que si le modèle réduit reste précis dans la bande de fréquence utile et pour les essais considérés.

% ==========================================================================
% COURS 05 -- PARTIE C : SCHÉMAS-BLOCS, BOUCLAGE ET PERTURBATIONS
% ==========================================================================

\section{Représentation par schéma-blocs}\label{sec:sa-schema-blocs}

Un schéma-blocs représente les relations entre les signaux d'un système. Il permet de conserver une lecture fonctionnelle tout en manipulant les fonctions de transfert.

\SASchemaElementsBase

\begin{itemize}
  \item Un \textbf{bloc} représente une relation $S(p)=H(p)E(p)$.
  \item Un \textbf{comparateur} réalise une somme algébrique.
  \item Un \textbf{point de prélèvement} duplique un signal sans le modifier.
  \item Les flèches imposent le sens de propagation de l'information.
\end{itemize}

\begin{SAattention}
Un schéma-blocs ne représente pas directement la circulation d'énergie. Les signaux peuvent être de natures et d'unités différentes. Chaque bloc doit donc être associé à une entrée, une sortie et une unité cohérentes.
\end{SAattention}

\subsection{Blocs en série, en parallèle et en boucle}

\TLReglesBlocs

Pour des blocs en série :
\[
H_{eq}(p)=\prod_i H_i(p).
\]

Pour des blocs en parallèle soumis à la même entrée :
\[
H_{eq}(p)=H_1(p)+H_2(p)
\]
avec les signes indiqués au sommateur.

Pour une rétroaction négative de chaîne directe $D(p)$ et de retour $R(p)$ :
\[
\boxed{H_{BF}(p)=\frac{D(p)}{1+D(p)R(p)}}.
\]
Pour une rétroaction positive, le signe du dénominateur devient négatif.

\subsection{Déplacement des sommateurs et points de prélèvement}

Lorsqu'un sommateur ou un point de prélèvement est déplacé de part et d'autre d'un bloc, il faut introduire le bloc ou son inverse afin de conserver les mêmes relations entre signaux.

\begin{SAmethode}[title={Simplifier un schéma-blocs}]
\begin{enumerate}
  \item Identifier les boucles internes et les réduire en premier.
  \item Regrouper les blocs en série ou en parallèle.
  \item Déplacer un sommateur ou un prélèvement uniquement si cela simplifie réellement la structure.
  \item Vérifier le résultat sur les dimensions et sur des cas limites simples.
\end{enumerate}
\end{SAmethode}

\section{Fonctions de transfert en boucle ouverte et fermée}\label{sec:sa-ftbo-ftbf}

Dans une boucle à retour unitaire :
\[
\varepsilon(p)=E(p)-S(p),\qquad S(p)=D(p)\varepsilon(p).
\]
On obtient :
\[
\frac{S(p)}{E(p)}=\frac{D(p)}{1+D(p)},
\qquad
\frac{\varepsilon(p)}{E(p)}=\frac{1}{1+D(p)}.
\]

Dans le cas général, la fonction de transfert en boucle ouverte est :
\[
L(p)=D(p)R(p).
\]
La fonction de transfert en boucle fermée est :
\[
H_{BF}(p)=\frac{D(p)}{1+L(p)}.
\]

Le polynôme $1+L(p)$ est l'équation caractéristique de la boucle fermée. Ses racines sont les pôles du système bouclé.

\begin{SAapplication}[title={Capteur non unitaire}]
Si la chaîne directe vaut $C(p)P(p)$ et le capteur $H(p)$ :
\[
\frac{S(p)}{E(p)}=\frac{C(p)P(p)}{1+C(p)P(p)H(p)}.
\]
La sortie ne suit pas nécessairement la consigne avec un gain statique égal à 1. L'étalonnage du capteur et les unités du comparateur doivent être pris en compte.
\end{SAapplication}

\section{Systèmes perturbés}\label{sec:sa-perturbations}

\TLSchemaPerturbe

Une perturbation peut intervenir à différents endroits de la chaîne. Son effet sur la sortie dépend de son point d'application.

\subsection{Théorème de superposition}

Pour un système linéaire à plusieurs entrées, la sortie totale est la somme des sorties dues à chaque entrée prise séparément :
\[
S(p)=H_E(p)E(p)+H_Q(p)Q(p).
\]

Pour déterminer $H_E$, on annule les autres entrées. Pour déterminer $H_Q$, on annule la consigne.

\begin{SAapplication}[title={Perturbation appliquée à l'entrée du procédé}]
Avec une chaîne directe $C(p)P(p)$, un retour unitaire et une perturbation $Q(p)$ ajoutée avant $P(p)$ :
\[
S(p)=\frac{C(p)P(p)}{1+C(p)P(p)}E(p)
+\frac{P(p)}{1+C(p)P(p)}Q(p).
\]
Un grand gain de boucle réduit l'influence de la perturbation dans la bande où $|C(j\omega)P(j\omega)|$ est grand.
\end{SAapplication}

\subsection{Sensibilité et rejet de perturbation}

On définit la fonction de sensibilité :
\[
S_e(p)=\frac{1}{1+L(p)},
\]
et la sensibilité complémentaire :
\[
T(p)=\frac{L(p)}{1+L(p)}.
\]
Elles vérifient :
\[
S_e(p)+T(p)=1.
\]

À basse fréquence, un grand gain de boucle rend $S_e$ faible : l'erreur et certaines perturbations sont rejetées. À haute fréquence, un gain trop grand amplifie le bruit de mesure et peut dégrader la robustesse.

\section{Précision statique et classe d'un système}\label{sec:sa-precision-statique}

Pour un retour unitaire :
\[
\varepsilon(p)=\frac{E(p)}{1+L(p)}.
\]
L'erreur statique est obtenue par le théorème de la valeur finale :
\[
\varepsilon_s=\lim_{p\to0}p\varepsilon(p).
\]

On écrit la boucle ouverte sous la forme :
\[
L(p)=\frac{K}{p^\alpha}\,L_0(p),\qquad L_0(0)=1.
\]
L'entier $\alpha$ est la \textbf{classe} du système.

\subsection{Erreur pour une consigne échelon}

Pour $E(p)=E_0/p$ :
\[
\varepsilon_s=\lim_{p\to0}\frac{E_0}{1+L(p)}.
\]
Ainsi :
\begin{itemize}
  \item classe 0 : $\varepsilon_s=E_0/(1+K)$ si $L(0)=K$ ;
  \item classe $\geq1$ : $\varepsilon_s=0$ si le système bouclé est stable.
\end{itemize}

\subsection{Erreur pour une consigne rampe}

Pour $E(p)=V_0/p^2$ :
\[
\varepsilon_v=\lim_{p\to0}\frac{V_0}{p(1+L(p))}.
\]
Ainsi :
\begin{itemize}
  \item classe 0 : erreur infinie ;
  \item classe 1 : erreur finie $V_0/K_v$ ;
  \item classe $\geq2$ : erreur nulle.
\end{itemize}

\begin{SAattention}
Ajouter un intégrateur améliore la précision à basse fréquence mais diminue généralement les marges de stabilité. La précision ne peut donc pas être améliorée indépendamment de la stabilité et de la rapidité.
\end{SAattention}

\section{Identifier un modèle de comportement}\label{sec:sa-identification}

La modélisation de comportement s'appuie sur des essais. Elle est utile lorsque les paramètres physiques sont inconnus ou lorsque le modèle de connaissance serait trop complexe.

\begin{SAmethode}[title={Démarche d'identification temporelle}]
\begin{enumerate}
  \item Placer le système dans un régime initial stable.
  \item Appliquer une entrée connue, souvent un échelon.
  \item Enregistrer l'entrée et la sortie avec une fréquence d'échantillonnage suffisante.
  \item Retirer les offsets et repérer un éventuel retard.
  \item Choisir une famille de modèle : premier ordre, second ordre, intégrateur, retard.
  \item Estimer les paramètres puis comparer la réponse simulée aux mesures.
\end{enumerate}
\end{SAmethode}

\subsection{Identification d'un premier ordre}

Pour un modèle $K/(1+\tau p)$ :
\[
K=\frac{\Delta s_\infty}{\Delta e},
\]
et $\tau$ est lu à $63{,}2\%$ de la variation finale. En présence d'un retard $T$, on utilise :
\[
H(p)=\frac{K e^{-Tp}}{1+\tau p}.
\]

\subsection{Identification d'un second ordre pseudo-périodique}

Pour une réponse présentant des dépassements :
\[
D_1=\exp\left(-\frac{\pi\xi}{\sqrt{1-\xi^2}}\right),
\]
d'où :
\[
\xi=\frac{-\ln D_1}{\sqrt{\pi^2+(\ln D_1)^2}}.
\]
La pseudo-période mesurée $T_a$ donne :
\[
\omega_a=\frac{2\pi}{T_a},\qquad
\omega_0=\frac{\omega_a}{\sqrt{1-\xi^2}}.
\]

\subsection{Validation}

Un modèle identifié doit être validé sur un essai différent de celui utilisé pour l'estimation. On compare notamment :
\begin{itemize}
  \item la valeur finale ;
  \item le temps de réponse ;
  \item les dépassements et la pseudo-période ;
  \item l'erreur quadratique entre mesure et simulation ;
  \item la cohérence des paramètres avec la physique du système.
\end{itemize}

% ==========================================================================
% COURS 05 -- PARTIE D : COMPORTEMENT FRÉQUENTIEL ET STABILITÉ
% ==========================================================================

\section{Réponse fréquentielle}\label{sec:sa-frequentiel}

Lorsqu'un SLCI stable est soumis à une entrée sinusoïdale de pulsation $\omega$, son régime permanent est sinusoïdal à la même pulsation. Seules l'amplitude et la phase sont modifiées.

Pour :
\[
e(t)=E_0\sin(\omega t),
\]
on obtient en régime permanent :
\[
s(t)=S_0\sin(\omega t+\varphi).
\]
La réponse fréquentielle est donnée par :
\[
H(j\omega)=\frac{S_0}{E_0}e^{j\varphi}.
\]

\begin{SAdefinition}
Le gain et la phase sont :
\[
G(\omega)=|H(j\omega)|,
\qquad
\varphi(\omega)=\arg H(j\omega).
\]
Le gain en décibels est :
\[
G_{dB}(\omega)=20\log_{10}|H(j\omega)|.
\]
\end{SAdefinition}

\subsection{Décade, octave et décibel}

Une décade correspond à un rapport de pulsations égal à 10. Une octave correspond à un rapport égal à 2.

Les propriétés logarithmiques transforment les produits en sommes :
\[
20\log_{10}|H_1H_2|=20\log_{10}|H_1|+20\log_{10}|H_2|.
\]
Ainsi, les diagrammes de Bode d'un produit de fonctions élémentaires s'obtiennent par addition des gains et des phases.

\section{Diagrammes de Bode élémentaires}\label{sec:sa-bode-elementaire}

\subsection{Gain pur}

Pour $H(p)=K$ :
\[
G_{dB}=20\log_{10}|K|,
\]
et la phase vaut $0^\circ$ si $K>0$, $-180^\circ$ ou $180^\circ$ si $K<0$.

\subsection{Intégrateur}

Pour $H(p)=K/p$ :
\[
G_{dB}=20\log_{10}|K|-20\log_{10}\omega,
\qquad
\varphi=-90^\circ.
\]
La pente du gain est $-20$ dB par décade.

\TLBodeIntegrateur[1]

\subsection{Premier ordre}

Pour :
\[
H(p)=\frac{K}{1+\tau p},
\]
la pulsation de cassure est :
\[
\omega_c=\frac1\tau.
\]

Les asymptotes du gain sont :
\[
G_{dB}\simeq20\log_{10}|K|\quad\text{pour }\omega\ll\omega_c,
\]
\[
G_{dB}\simeq20\log_{10}|K|-20\log_{10}\left(\frac{\omega}{\omega_c}\right)
\quad\text{pour }\omega\gg\omega_c.
\]
À la cassure, le gain réel est inférieur de 3 dB à l'asymptote et la phase vaut $-45^\circ$.

\TLBodePremierOrdre[1][1]

\subsection{Second ordre}

Pour :
\[
H(p)=\frac{K\omega_0^2}{p^2+2\xi\omega_0p+\omega_0^2},
\]
le gain présente une pente de $-40$ dB par décade à haute fréquence et la phase évolue de $0^\circ$ vers $-180^\circ$.

Pour $\xi<1/\sqrt2$, un pic de résonance apparaît. Sa pulsation est :
\[
\omega_r=\omega_0\sqrt{1-2\xi^2}.
\]

\TLBodeSecondOrdre[1][0.7][1]

\begin{SAmethode}[title={Tracer un diagramme de Bode asymptotique}]
\begin{enumerate}
  \item Mettre la fonction de transfert sous forme factorisée canonique.
  \item Repérer le gain, les intégrateurs, les dérivateurs et les pulsations de cassure.
  \item Tracer le gain initial puis modifier la pente à chaque cassure.
  \item Additionner les contributions de phase.
  \item Placer quelques points réels remarquables autour des cassures.
\end{enumerate}
\end{SAmethode}

\section{Stabilité en boucle fermée}\label{sec:sa-stabilite}

\subsection{Critère sur les pôles}

Un SLCI continu est asymptotiquement stable si tous les pôles de sa fonction de transfert sont à partie réelle strictement négative.

\begin{itemize}
  \item Pôle réel négatif : mode exponentiel décroissant.
  \item Paire complexe à partie réelle négative : oscillations amorties.
  \item Pôle à partie réelle positive : divergence.
  \item Pôle simple sur l'axe imaginaire : stabilité marginale selon le contexte.
\end{itemize}

\subsection{Point critique}

Pour une rétroaction négative, l'équation caractéristique est :
\[
1+L(p)=0.
\]
La limite d'instabilité correspond donc au point :
\[
L(j\omega)=-1,
\]
soit un gain de $0$ dB et une phase de $-180^\circ$.

\begin{SAdefinition}
La \textbf{marge de phase} est la distance angulaire à $-180^\circ$ lorsque le gain de boucle vaut 0 dB :
\[
M_\varphi=180^\circ+\arg L(j\omega_{0dB}).
\]
La \textbf{marge de gain} est l'opposé du gain en dB lorsque la phase vaut $-180^\circ$ :
\[
M_G=-G_{dB}(\omega_{-180^\circ}).
\]
\end{SAdefinition}

Des marges positives sont nécessaires à la stabilité, mais une marge trop faible rend le système fragile aux incertitudes et aux retards non modélisés.

\begin{SAattention}
La stabilité se juge sur la boucle ouverte $L(p)$, mais concerne les pôles de la boucle fermée. Il ne faut pas confondre la stabilité propre d'un bloc et la stabilité de l'asservissement complet.
\end{SAattention}

\subsection{Lien entre bande passante et rapidité}

Une augmentation de la bande passante de la boucle fermée rend généralement la réponse plus rapide. Toutefois, une bande passante trop élevée augmente la sensibilité au bruit, aux dynamiques négligées et aux retards.

Une approximation courante pour un comportement dominé par un premier ordre est :
\[
t_{5\%}\approx\frac{3}{\omega_B},
\]
où $\omega_B$ est une pulsation représentative de la bande passante. Cette relation reste indicative et dépend de la forme réelle de la réponse.

\section{Bilan de la démarche d'étude}\label{sec:sa-demarche}

\begin{SAmethode}[title={Étudier un système asservi}]
\begin{enumerate}
  \item \textbf{Définir le système :} consigne, sortie, commande, perturbations, capteurs et unités.
  \item \textbf{Construire le modèle :} lois physiques ou identification expérimentale.
  \item \textbf{Établir le schéma-blocs :} fonctions de transfert et signes des sommateurs.
  \item \textbf{Calculer FTBO et FTBF :} simplifier sans perdre le sens physique.
  \item \textbf{Vérifier la stabilité :} pôles ou marges fréquentielles.
  \item \textbf{Évaluer les performances :} précision, rapidité, dépassement et rejet des perturbations.
  \item \textbf{Comparer au cahier des charges :} conclure quantitativement.
  \item \textbf{Valider le modèle :} confrontation aux mesures et analyse des écarts.
\end{enumerate}
\end{SAmethode}

\begin{center}
\begin{tabularx}{\linewidth}{>{\bfseries}p{.19\linewidth}X X}
\toprule
Performance & Domaine temporel & Domaine fréquentiel \\
\midrule
Stabilité & pôles, oscillations, divergence & marges de phase et de gain \\
Rapidité & temps de réponse, temps de montée & bande passante, pulsation de coupure \\
Précision & erreur statique, erreur de traînage & gain de boucle à basse fréquence \\
Dépassement & maximum de la réponse indicielle & résonance, amortissement \\
Robustesse & sensibilité aux paramètres & marges, sensibilité, bruit \\
\bottomrule
\end{tabularx}
\end{center}

\section{Sources}\label{sec:sa-sources}

Ce chapitre reprend et réorganise les contenus des cours d'introduction et de modélisation des systèmes asservis. Les méthodes et notations sont harmonisées avec le framework SII du dépôt et avec les usages classiques de l'automatique continue en CPGE et en ATS.

% ==========================================================================
% COURS 05 -- EXERCICES
% ==========================================================================

\section{Exercices du chapitre}\label{sec:sa-exercices}

\begin{TLExercise}{Exercice 1 -- Étuve thermique}
Une étuve est commandée par un gradateur. La température est mesurée par une sonde Pt100 et comparée à une consigne.
\begin{enumerate}
  \item Identifier la consigne, la sortie, la commande et deux perturbations possibles.
  \item Proposer un schéma-blocs fonctionnel.
  \item Expliquer pourquoi une commande en boucle ouverte serait insuffisante.
  \item Préciser s'il s'agit principalement d'un asservissement ou d'une régulation.
\end{enumerate}
\end{TLExercise}

\begin{TLExercise}{Exercice 2 -- Fauteuil motorisé}
Le dossier d'un fauteuil est déplacé par un motoréducteur. Un potentiomètre mesure l'angle réel.
\begin{enumerate}
  \item Compléter la chaîne : comparateur, correcteur, motoréducteur, mécanisme, capteur.
  \item Donner les unités des signaux au niveau du comparateur.
  \item Citer un défaut de capteur susceptible de créer une erreur permanente.
  \item Expliquer l'effet d'une augmentation du gain du correcteur sur précision, rapidité et stabilité.
\end{enumerate}
\end{TLExercise}

\begin{TLExercise}{Exercice 3 -- Fonction de transfert d'un circuit RC}
Un circuit vérifie :
\[
RC\frac{du_s(t)}{dt}+u_s(t)=u_e(t).
\]
\begin{enumerate}
  \item Déterminer $H(p)=U_s(p)/U_e(p)$ pour des conditions initiales nulles.
  \item Identifier le gain statique et la constante de temps.
  \item Donner la réponse à un échelon d'amplitude $E_0$.
  \item Calculer le temps de réponse à 5\%.
\end{enumerate}
\end{TLExercise}

\begin{TLExercise}{Exercice 4 -- Moteur à courant continu}
On considère :
\[
u=Ri+L\dot{i}+K_e\omega,
\qquad
J\dot{\omega}+f\omega=K_t i,
\]
et un couple résistant nul.
\begin{enumerate}
  \item Transformer les équations dans le domaine de Laplace.
  \item Éliminer $I(p)$ et déterminer $\Omega(p)/U(p)$.
  \item Donner l'ordre du modèle et identifier ses paramètres physiques.
  \item Proposer une approximation du premier ordre lorsque l'inductance est négligeable.
\end{enumerate}
\end{TLExercise}

\begin{TLExercise}{Exercice 5 -- Identification d'un premier ordre}
Un essai indiciel de $2$ V produit une variation finale de sortie de $18$ mm. La sortie atteint $11{,}4$ mm après $0{,}42$ s.
\begin{enumerate}
  \item Identifier le gain statique $K$.
  \item Identifier la constante de temps $\tau$.
  \item Estimer le temps de réponse à 5\%.
  \item Écrire la fonction de transfert et la réponse temporelle.
\end{enumerate}
\end{TLExercise}

\begin{TLExercise}{Exercice 6 -- Second ordre pseudo-périodique}
La réponse à un échelon unitaire tend vers $2$. Le premier maximum vaut $2{,}5$ et deux maxima successifs sont séparés de $0{,}8$ s.
\begin{enumerate}
  \item Calculer le dépassement relatif.
  \item En déduire le coefficient d'amortissement $\xi$.
  \item Déterminer la pulsation amortie puis la pulsation propre $\omega_0$.
  \item Écrire une fonction de transfert canonique compatible avec l'essai.
\end{enumerate}
\end{TLExercise}

\begin{TLExercise}{Exercice 7 -- Réduction d'un schéma-blocs}
Une chaîne directe est constituée de trois blocs $K$, $1/(1+\tau p)$ et $1/p$ en série. Le retour est unitaire et négatif.
\begin{enumerate}
  \item Déterminer la FTBO.
  \item Déterminer la FTBF.
  \item Donner l'équation caractéristique.
  \item Préciser la classe du système et son erreur statique pour une consigne échelon.
\end{enumerate}
\end{TLExercise}

\begin{TLExercise}{Exercice 8 -- Perturbation de couple}
Un moteur est commandé par un correcteur $C(p)$ et un procédé $P(p)$. Une perturbation $Q(p)$ s'ajoute à l'entrée du procédé. Le retour est unitaire.
\begin{enumerate}
  \item Déterminer $S(p)/E(p)$ en annulant $Q(p)$.
  \item Déterminer $S(p)/Q(p)$ en annulant $E(p)$.
  \item Montrer comment un grand gain de boucle améliore le rejet de la perturbation.
  \item Expliquer pourquoi ce gain ne peut pas être augmenté sans limite.
\end{enumerate}
\end{TLExercise}

\begin{TLExercise}{Exercice 9 -- Erreurs statiques}
On considère un retour unitaire de boucle ouverte :
\[
L(p)=\frac{20}{p(1+0{,}2p)}.
\]
\begin{enumerate}
  \item Déterminer la classe du système.
  \item Calculer l'erreur statique pour un échelon d'amplitude 3.
  \item Calculer l'erreur de traînage pour une rampe de pente $0{,}5$.
  \item Conclure sur la précision du système.
\end{enumerate}
\end{TLExercise}

\begin{TLExercise}{Exercice 10 -- Diagramme de Bode}
On considère :
\[
H(p)=\frac{50}{p(1+0{,}1p)(1+p)}.
\]
\begin{enumerate}
  \item Mettre la fonction sous forme canonique et repérer les pulsations de cassure.
  \item Tracer les asymptotes du gain.
  \item Tracer l'évolution approchée de la phase.
  \item Donner les pentes successives du gain en dB par décade.
\end{enumerate}
\end{TLExercise}

\begin{TLExercise}{Exercice 11 -- Marges de stabilité}
Pour une boucle ouverte, la pulsation de coupure à 0 dB vaut $8$ rad/s et la phase correspondante vaut $-145^\circ$. À la pulsation où la phase vaut $-180^\circ$, le gain vaut $-9$ dB.
\begin{enumerate}
  \item Calculer la marge de phase.
  \item Calculer la marge de gain.
  \item Commenter la robustesse qualitative de la boucle.
  \item Indiquer l'effet probable d'un retard supplémentaire.
\end{enumerate}
\end{TLExercise}

\begin{TLExercise}{Exercice 12 -- Synthèse complète}
Un axe linéaire doit suivre une consigne de position. Le cahier des charges impose : erreur statique nulle sur un échelon, temps de réponse à 5\% inférieur à $0{,}8$ s et dépassement inférieur à $10\%$.
\begin{enumerate}
  \item Proposer le protocole d'essai permettant de mesurer ces trois performances.
  \item Identifier les informations nécessaires pour construire un modèle.
  \item Établir une démarche de calcul de la FTBF à partir du schéma-blocs.
  \item Indiquer les grandeurs à lire dans les domaines temporel et fréquentiel.
  \item Expliquer les compromis à respecter lors du choix du correcteur.
\end{enumerate}
\end{TLExercise}

\end{document}
